\begin{table}[htbp]\centering\small
\caption{Balance, RFSP y saldo de la deuda}\label{cua:deuda}
\begin{tabular}{lrrr}
\toprule
Concepto & 2026 aprob. & 2027 & 2027 (\% PIB) \\
\midrule
Balance presupuestario (mdp) & $-$1.393.770,6 & $-$1.358.558,8 & $-$3,4 \\
Superávit económico primario (mdp) & 178.802,6 & 216.891,3 & 0,6 \\
RFSP (mdp) & $-$1.587.349,9 & $-$1.555.655,8 & $-$3,9 \\
SHRFSP (mdp) & --- & 21.665.995,8 & 55,0 \\
\quad interno (\% PIB) & --- & --- & 43,6 \\
\quad externo (\% PIB) & --- & --- & 11,4 \\
Endeudamiento neto interno autorizado (mdp) & --- & 1.700.000,0 & --- \\
Endeudamiento neto externo autorizado (mdd) & --- & 13.500,0 & --- \\
\bottomrule
\end{tabular}
\\[2pt]\begin{minipage}{0.92\textwidth}\scriptsize Fuentes: CGPE 2027, edición SHCP, Anexos II.6 (p. 67) y III.2 (p. 47); artículo 2o. de la ILIF 2027. \textbf{El superávit primario se toma del Anexo II.6 (0,6~\% del PIB) y no del cuadro resumen de la página 6 de esa edición, que imprime 0,5 y contradice a su propio anexo.} \textbf{Techo de endeudamiento y déficit son objetos distintos}: los dos últimos renglones autorizan, no proyectan.\end{minipage}
\end{table}
